\newcommand{\DoNotLoadEpstopdf}{}
\documentclass[11pt,reqno]{amsart}
\usepackage[T1]{fontenc}
\usepackage{lmodern,amsmath,amssymb,amsthm,mathtools,mathrsfs,microtype}
\usepackage[margin=1.1in]{geometry}
\usepackage{xr-hyper}
\usepackage[hidelinks]{hyperref}
\externaldocument[supp-]{build/supplement}[supplement.pdf]
\hypersetup{pdftitle={Collision records, arithmetic, and raw local inversion in a triangular Lorentz gas},pdfauthor={Qian Qi}}
\setlength{\emergencystretch}{2em}
\allowdisplaybreaks[2]
\numberwithin{equation}{section}
\newtheorem{theorem}{Theorem}[section]
\newtheorem{lemma}[theorem]{Lemma}
\newtheorem{proposition}[theorem]{Proposition}
\newtheorem{corollary}[theorem]{Corollary}
\theoremstyle{definition}
\newtheorem{definition}[theorem]{Definition}
\newtheorem{inputstatement}[theorem]{Input statement}
\theoremstyle{remark}\newtheorem{remark}[theorem]{Remark}
\newcommand{\R}{\mathbb R}\newcommand{\Z}{\mathbb Z}
\newcommand{\one}{\mathbf1}\newcommand{\dd}{\,\mathrm d}
\newcommand{\esssup}{\operatorname*{ess\,sup}}
\newcommand{\diver}{\operatorname{div}}
\newcommand{\TV}{\operatorname{TV}}
\title[Collision records and raw local inversion]{Collision records, arithmetic, and raw local inversion in a triangular Lorentz gas}
\author{Qian Qi}
\date{October 11, 2026. A2-DYN, revision 82}
\subjclass[2020]{37D50, 37A50, 60F05, 28A75}
\keywords{Sinai billiards, exact collision records, local limits, joint endpoint localization, mixed singular insertions, optical drift, exact-label coarea flux}
\begin{document}
\raggedbottom
\begin{abstract}
We retain the pointwise local limit problem for the full displacement,
collision-count, occupation and roof record of actual returns in a
triangular Lorentz gas. Joint localization of the two endpoint momenta
and one internal incidence yields a superlinear gain in both widths.
The proof distinguishes a small initial vector from a bounded
multiplier, and treats both chronological cases at every continued
arithmetic peak. A physical cap on the finite-count remainder permits
all inverse-incidence scales to be integrated. The resulting complete
optical source has a critical-width gain greater than one, actual
untruncated fractional inverse-gradient moments, and first
inverse-gradient bounds uniform in a fixed positive lower cutoff.
Applied on common exact-label coarea fibers, these estimates replace
the earlier optical and grazing-exit bound $Cr/a$ by $Cr$, and pay
the same uniform compression cost at nongrazing exits. The complete
critical initial defect and outgoing nongrazing flux remain the two
roof-height estimates needed for the original pointwise theorem;
no exchange of critical-cutoff and collision limits is inferred.
\end{abstract}
\maketitle

\section{The unchanged physical problem}
Put $R_-=9/20$ and $R_+=47/100$, and let $R\in I=[R_-,R_+]$. Circular scatterers have centers in
$\Lambda=\Z(1,0)+\Z(1/2,\sqrt3/2)$. Write $T_R$ for the regular
collision map and $(\alpha,p)\in(\R/2\pi\Z)\times(-1,1)$ for its
coordinates; $p$ is the outgoing tangential velocity and
$\dd\nu=\dd\alpha\,\dd p/(4\pi)$. Physical flights satisfy
$\tau_-\le\tau_R\le\tau_+$, where $\tau_-=3/50$ and $\tau_+$
is the fixed finite-horizon bound.

The section is specified here to fix every event in the argument.
Put $h_0=1/200$. Let $Y_R$ be the union of the four open coordinate
squares of half-side $h_0$ centered at
\[
 (0,0),\quad(\pi/6,-1/2),\quad
 (-\theta_R,\sin\theta_R),\quad
 (\pi/3-\theta_R,\sin\theta_R).
\]
Here $\theta_R=2\arctan q_R$, and $q_R$ is the unique root in
$(49/100,53/100)$ of
\[
 8Rq(1+q^2)-4q(1-q^2)+\sqrt3(1-6q^2+q^4)=0.
\]
Set
\[
 U=\{|\alpha-\pi/6|<3/50,\ |p|<3/20\},\qquad
 Y_R^*=Y_R\cup U,\qquad c_*=\frac{91}{10000\pi}.
\]
The rectangles are disjoint and $\nu(Y_R^*)=c_*$. We always use
$\nu_R^*=c_*^{-1}\nu|_{Y_R^*}$, not a renormalization of a smaller
source.

At $m$ flights, let $L=L_{m,R}$ be their total length,
$K^{\rm c}_{m,R}$ their lattice displacement, and
$A_{m,R}=\sum_{j=0}^{m-1}\one_{Y_R^*}\circ T_R^j$. The event is
\begin{equation}\label{eq:v76-event}
 \Pi_{n,k,m,R}=Y_R^*\cap\{K^{\rm c}_{m,R}=k,\ A_{m,R}=n\}
                         \cap T_R^{-m}Y_R^*.
\end{equation}
The initial visit is counted; the terminal visit is tested separately.
On this event the actual $n$th return has collision count $m$ and
roof $L$. Denote its density by $p_{n,R}(k,m,t)$.

Write $c_j=\sqrt{1-p_j^2}$ and let $\Delta_j$ be the least clearance
of flight $j$ from any third disk, excluding its two incident disks.
Keep the original smooth function $0\le\Theta\le1$, zero on
$(-\infty,1]$ and one on $[2,\infty)$. With $q_0=47/53$,
$d_j=\min(j,m-j)$ and $d'_j=\min(j,m-1-j)$, set
\begin{align}
 H^\varepsilon&=
 \prod_{j=0}^{m}\Theta\!\left(\frac{c_j}{\varepsilon q_0^{d_j/4}}\right)
 \prod_{j=0}^{m-1}\Theta\!\left(\frac{\Delta_j}{\varepsilon q_0^{d'_j/4}}\right),
 \qquad W^\varepsilon=1-H^\varepsilon,\label{eq:v76-guard}\\
 b^\varepsilon(t)\dd t
   &=L_\#(W^\varepsilon\one_{\Pi_{n,k,m,R}}\nu_R^*).
 \label{eq:v76-source}
\end{align}
In particular first incidence and first clearance outside any selected
critical neighborhoods are still present. No section-distance factor
has been inserted into this guard.

\subsection{The arithmetic reference and the target}
For $y=(k_1,k_2,t,n)$ put
\[
 \mu_R=(0,0,\bar\tau_R,c_*),\quad
 Z=(y-m\mu_R)/\sqrt m,\quad
 \mathcal C_{m,R,n,k}(M)=\{t:|Z|\le M\}.
\]
We retain the finite spectral-peak reference, including peaks close to
but not on the unit circle. To specify its meaning, the collision
record observable is $(\kappa_1,\kappa_2,\tau,\one_{Y_R^*})$.
Let $\lambda_j(R,\vartheta)$ and $\mathsf P_j(R,\vartheta)$ be its
continued simple twisted eigenvalue and projection in each of the
fixed peak charts, and let $\chi_j$ be their frequency partition.
At the real maximum $a_j$ of $\log|\lambda_j|$ put
\[
 z_j=\lambda_j(R,a_j),\quad
 \mu_j=i^{-1}\nabla\log\lambda_j(R,a_j),\quad
 H_j=-\nabla^2\log\lambda_j(R,a_j),
\]
with $\operatorname{Re}H_j\ge aI_4$ for a common $a>0$. If
$\eta_R=\one_{Y_R^*}$ and $\ell$ denotes mass evaluation, put
$B_j=\chi_j(R,a_j)\ell(M_{\eta_R}\mathsf P_j(R,a_j)(\eta_R\nu))$.
Define
\begin{align}
 G_H(z)&=(2\pi)^{-4}\int_{\R^4}e^{-iv\cdot z-v^{\mathsf T}Hv/2}\dd v,
 \notag\\
 \mathcal L_{m,R}(y)&=\operatorname{Re}\left[
 c_*^{-1}\sum_j e^{-ia_j\cdot y}z_j^m B_j
      G_{H_j}\left((y-m\mu_j)/\sqrt m\right)\right].
 \label{eq:v76-reference}
\end{align}
Only the charts active at $R$ are summed. This is the same transition
reference as in the retained spectral theorem. The existence and
uniform estimates for these scalar data belong to that analytic
input, not to the transport argument below. No frequency branch or
arithmetic zero class is deleted.

The governing assertion remains
\begin{equation}\label{eq:v76-target}
 \sup_{R,\,1\le n\le m,\,k\in\Z^2}
 \esssup_{t\in\mathcal C_{m,R,n,k}(M)}
 |m^2p_{n,R}(k,m,t)-\mathcal L_{m,R}(k_1,k_2,t,n)|\longrightarrow0.
\end{equation}
Let $\mathcal E_M$ denote this error's collision limsup. An empty
central target contributes zero. The order of the uniformity is
$\sup_R\esssup_t$, so exceptional density sets may depend on $R$.

\subsection{The analytic inputs in the form used here}
\begin{inputstatement}\label{input:v76-analytic}
The unchanged source \eqref{eq:v76-source} has a nonnegative,
compactly supported $BV$ density at every fixed count. For every
fixed $\varepsilon,\delta>0$, with
$A_\delta f(t)=(2\delta)^{-1}\int_{t-\delta}^{t+\delta}f(s)\dd s$,
\begin{equation}\label{eq:v76-average}
 \limsup_m\sup_{R,n,k}m^2\|A_\delta b^\varepsilon\|_\infty
                                      \le C\varepsilon^{1/16}.
\end{equation}
The constant is independent of the previously fixed width. For large
fixed $B$ and $\varepsilon(B)=A_0B^{-1/12}$,
\begin{equation}\label{eq:v76-remainder}
 m^2p-\mathcal L=m^2b^{\varepsilon(B)}+e_B,\qquad
 \limsup_m\sup_{R,n,k}\esssup_{\mathcal C(M)}|e_B|
                                     \le C_MB^{-1/192}.
\end{equation}
All auxiliary spectral bands in these statements are fixed before
the collision limit and enlarged only afterwards.
\end{inputstatement}
These are the complete-source tame theorem, the fixed-width
complete-source bound, and the positive scalar remainder theorem in
the technical supplement. Their proofs and historical hypotheses are
retained there. The fixed-count source uses threshold integration
\cite{CMInt}, definable monotonicity \cite{Coste}, and the
distributional characterization of $BV$ \cite{AFP}; it does not
assume definability of the original smooth guard. They are restated as inputs to distinguish them from
the new geometric and transport proofs. No long-time bound on a
fixed-count variation constant is part of this statement.



\subsection{Precise inherited inputs}\label{sec:v79-input-map}
The fixed-width estimate is Supplement Theorem~\ref{supp-thm:v49-physical-remainder-local},
in particular Supplement equation~\eqref{supp-eq:v49-local-limsup}.
It uses Supplement Theorem~\ref{supp-thm:v49-physical-marked-local}, on the
original exact displacement and occupation coefficient. For an interval
of length $2r$, division by $2r$ gives \eqref{eq:v76-average}. That
proof fixes the auxiliary band, takes the collision limit, and only
then enlarges the auxiliary band. Its leading constant is independent
of the previously fixed interval length. The extension from small
margins to a fixed remaining compact range is stated there explicitly.
At auxiliary parameter one, $\Theta(c_0)=0$ since $c_0\le1$, so $W^1=1$
and $b^1=p$. This supplies the unweighted fixed-width estimate without
asserting a raw pointwise bound. The positive remainder is the inherited
scalar input \eqref{eq:v76-remainder}, stated as Supplement
Theorem~\ref{supp-thm:v51-positive-error}; no surface-current estimate is
included in any of these volume-source inputs.

\subsection{The complete-source estimate proved here}
The physical endpoint variables are
$q=(\alpha_0,\alpha_m)$, $g=\nabla_qF=R(-p_0,p_m)$ and
$K=D_q^2F$. The retained canonical field satisfies $XF=1$,
$X\arg g=0$ and
$|g(\Phi_sx)|^2\ge |g(x)|^2+2\lambda_\circ s$ up to its true
lifetime. Its physical Jacobian factors into an inverse-twist ratio
and an angular Hessian ratio. The latter and the original guard
already have complete-source estimates.

The remaining optical drift is bounded by
$C\omega I_m/|g|$, where
\[
 I_m=1+\sum_{j=1}^{m-1}q_0^{d_j}/c_j.
\]
The preceding singular-insertion argument controlled $I_m$ but
replaced $|g|^{-1}$ by $a^{-1}$ on initial $|g|\ge a$.
The new argument keeps both components of $g$ in the spectral
pairing. With $\alpha=129/128$, it proves, for the complete original
source and each fixed roof width,
\[
 u^b(t)\,dt=L_\#(I_m\one_{\{|g|<b\}}\one_\Pi\nu_R^*),\qquad
 \limsup_m\sup_{R,n,k}m^2\|A_ru^b\|_\infty\le Cb^\alpha.
\]

Theorem~\ref{thm:v81-gradient-insertions} integrates this bound with
its finite-count remainder. It gives an actual untruncated mixed
insertion for powers $|g|^{-\eta}$ with $0<\eta<1/2$, and a bound
for $|g|^{-1}\one_{\{|g|\ge a\}}$ whose leading constant is
independent of the previously fixed $a>0$.
Theorem~\ref{thm:v81-optical-height} then proves a complete
$Cr$ optical and grazing starting-mass estimate on that initial
source. Nongrazing exits retain their exact outgoing flux; only
the paid pullback cost is improved. These statements keep the
original target \eqref{eq:v76-target}, rather than replacing it by
a local law for a different record.

\input{core/174_canonical_endpoint_momentum_flow}
\input{core/177_all_contact_motion_and_guard}
\input{core/180_uniform_cubic_reduced_action}
\input{core/181_separated_physical_volume_budget}
\input{core/182_optical_sensitivity_and_exit_flux}
\input{core/183_homogeneous_incidence_localization}
\input{core/184_singular_optical_transport}
\input{core/185_common_coarea_and_remaining_boundary}
\input{core/186_joint_endpoint_incidence_localization}
\input{core/187_mixed_critical_volume_insertions}
\input{core/188_uniform_critical_truncation_transport}

\section{The remaining scalar estimates and their consequences}
Equation~\eqref{eq:v81-endpoint-budget} is the revised error budget
for the original pointwise assertion. It retains the complete
critical initial defect and the nongrazing outgoing flux. The
uniformity of the new estimate removes the former restriction
$r/a\to0$ from the paid drift; it does not remove those two sources.
In particular the inverse-Jacobian cost has been estimated rather
than erased from the first-exit formula.

The quantitative distinction between a truncated and an untruncated
first inverse gradient is explicit in
\eqref{eq:v81-truncated-finite}. The actual fractional insertion
\eqref{eq:v81-fractional-finite} integrates both singular parameters
before the collision limit, whereas the first inverse power still
requires its positive cutoff to be fixed. The full pointwise theorem
is not asserted until the two remaining roof-height bounds are
proved in the stated order. All arithmetic peaks and zero classes
remain in the same reference.

After scalar closure, Supplement
Theorem~\ref{supp-thm:v73-scalar-path-equality} and the positive path
remainder give the same-roof path numerator. Normalization is only
on positive arithmetic reference classes. Collision-to-return clock
transfer and forward likelihood retain their separate inherited
hypotheses. The complete preceding article and every inherited proof
are retained in the supplement.

\input{core/189_dyadic_critical_source_criterion}

\input{journal-references-v81}
\end{document}
